\documentclass[Japanese]{dicomopapers}
%\documentclass[Japanese,noauthor]{dicomopapers}
\usepackage[dvipdfmx]{graphicx}
\usepackage{latexsym}
\usepackage{amsmath,amssymb,amsfonts}
\usepackage{otf}
\usepackage[]{siunitx}
\usepackage{booktabs}
\usepackage{tabularx}

\def\Underline{\setbox0\hbox\bgroup\let\\\endUnderline}
\def\endUnderline{\vphantom{y}\egroup\smash{\underline{\box0}}\\}
\def\|{\verb|}

%bibのエラー対応
\makeatletter
\newcommand\newblock{\hskip .11em\@plus.33em\@minus.07em}
\makeatother

\begin{document}

% 和文表題
\title{生活支援ロボットの安全な作業遂行に向けた\\力触覚情報に基づく接触状態の認識}

% 英文表題
\etitle{Haptic-Based Contact State Recognition for Safe Task Execution\\ in Home Service Robots}

\affiliate{UTokyoEEIS}{東京大学 大学院工学系研究科\\Graduate School of Engineering, The University of Tokyo}


\author{{\UTF{9AD9}}橋 和希}{KAZUKI TAKAHASHI}{UTokyoEEIS}[kazuki.takahashi@akg.t.u-tokyo.ac.jp]
\author{村上 弘晃}{HIROAKI MURAKAMI}{UTokyoEEIS}[murakami@akg.t.u-tokyo.ac.jp]
\author{佐久間 亮太}{RYOTA SAKUMA}{UTokyoEEIS}[sakuma@akg.t.u-tokyo.ac.jp]
\author{{\UTF{2EF2}\UTF{FA11}}允啓}{MITSUHIRO KAMEZAKI}{UTokyoEEIS}[kamezaki@akg.t.u-tokyo.ac.jp]
\author{川原 圭博}{YOSHIHIRO KAWAHARA}{UTokyoEEIS}[kawahara@akg.t.u-tokyo.ac.jp]

\begin{abstract}
Physical AIの発展に伴い，複雑な家庭内タスクを自律的に遂行する生活支援ロボットの実用化が期待されている．
現在，ロボットの環境認識や行動生成は視覚ベースのモデルが広く用いられているが，視覚のみでは対象物の重さや接触状態を十分に把握することは難しい．そのため，視覚情報を補完する手段として触覚情報の活用が進んでおり，特に環境との相互作用を高時間分解能で捉える力覚データの有用性が注目されている．
しかし，動的な力触覚時系列データから物理的状態や現象に対する意味を直接抽出し，ロボットの行動決定に活用する手法は十分に確立されていない．
本研究では，6軸力覚センサから得られる時系列データを，対照学習を用いて自然言語の潜在空間へマッピングすることで，ロボットが力学的現象を意味論的に解釈可能にする手法を提案する．
未知物体を用いた特徴量評価の結果，提案モデルは滑りや接地，物体の重量といった物理状態を高い精度で識別できることが確認された．さらに，「接地の手応えに合わせた把持の解除」や「重量ベースの物体選別」といった実用的なユースケースにおいて，個別の閾値設計を必要とせず，視覚ベースラインを上回る成功率を達成した．
力学的応答を汎用的な意味空間で解釈する本アプローチは，ロボットが実世界の物理的リスクを判断する際の手がかりを与える．これにより，人間の生活空間に溶け込み，安全に作業を遂行する生活支援ロボットの実現に貢献することが期待される．
\end{abstract}


\maketitle

\section{はじめに}
Physical AIの発展に伴い，家事などの複雑な家庭内タスクを自律的に遂行するロボットの実現が期待されている．
人の居住空間で活動するロボットには，単に対象物を把持して移動させるだけでなく，物体や環境との相互作用を理解し，状況に応じて安全かつ適切に振る舞う能力が求められる．
例えば，物体を机上に置く際には接地の瞬間を検知して適切に把持を解除する必要があり，飲食後の片付けでは外観が似た容器の中から空のものを重さに基づいて識別する必要がある．
これらの判断には，視覚情報のみでは捉えにくい接触状態や重さを把握するための触覚情報が重要となる．

しかし，従来の触覚センシングおよびその解釈手法は，ミリ秒単位で変化する動的な物理イベントの意味理解という要求を十分に満たしていない．
GelSight~\cite{GelSight}に代表される視触覚センサを用いたアプローチ~\cite{BindingTouch, CLTP, octo}の多くは，接触面の微細形状を高精度に観測できる一方で，画像処理に伴う計算負荷やサンプリング遅延により，衝撃や初期滑りのような高速な動的現象への即応には課題が残る．
また，触覚情報を入力とする行動生成モデル\cite{VTLA, tactilevla, ForceVLA}も存在するが，その多くはタスク成功率の向上や教師データの模倣のみを目的としたEnd-to-End学習に留まっており，状況を分析して生じているリスクを自律的に判断・応用するための抽象的な意味空間の獲得には至っていない．

そこで本研究では，6軸力覚センサから得られる高周波な時系列データを自然言語の潜在空間へ統合し，ロボットが物理現象を意味論的に解釈可能な力触覚エンコーダモデルを提案する．
具体的には，力触覚時系列の特徴量を言語モデルのベクトル空間と対照学習により対応付けることで，力覚の微細な変動を「滑り」や「接地」といった言語的意味として扱えるようにする（図~\ref{fig:overview}）．
力覚データは画像ベースの視触覚データに比べて軽量であり，\qty{100}{\Hz}以上の時間分解能で接触イベントを捉えられるため，高速な物理変化をリアルタイムに解釈するうえで有効である．
さらに，力学的応答を言語空間に射影することで，指先の局所的な接触状態だけでなく，「把持中の物体が机に触れた」といった物体を介した環境との相互作用を解釈できることが期待される．

本研究では，ロボットがユーザの要望に応じて何かを運び届けるユースケースを想定し，シミュレーション環境を用いたPick and Placeタスクによるデータ収集，モデル学習，評価を行った．
実験では，物理状態分類，特徴空間の可視化，および接触状態の理解と物体重量の推定を必要とする2種類のアプリケーション評価を通じて，提案手法の有効性を検証した．
その結果，提案モデルは物理状態分類において高い性能を示し，視覚ベースの手法と比較してアプリケーション成功率も向上することを確認した．

本研究の貢献は以下の通りである．
\begin{itemize}
    \item 力覚センサから得られる時系列データが，接触状態や物体重量の理解に有効な知覚情報源となることを示した．
    \item 力触覚時系列データを言語空間に対照学習で射影することで，動的な物理イベントを言語的に解釈可能な表現として獲得できることを示した．
    \item 力触覚データセットの自動生成パイプラインを構築し，複数のモデルアーキテクチャの比較を通じて，力触覚と言語を結びつけた物理状態理解の探索基盤を整備した．
\end{itemize}

\begin{figure*}[t]
    \centering
    \includegraphics[width=\linewidth]{fig/DICOMO_overview.pdf}
    \caption{提案する力触覚解釈システムの概要．(a) データ収集と自動アノテーション：シミュレータを用いたPick-and-Place動作からのマルチモーダルデータ収集，および時間窓分割による言語文の自動付与．(b) 対照学習によるアライメント：力触覚時系列と言語のペアデータを用いた対照学習，およびS-BERT~\cite{SBERT}の言語空間へ力触覚データを射影するエンコーダの構築．(c) アプリケーション評価：構築したエンコーダを用いた「接地の手応えに合わせた把持の解除」および「重量ベースの物体選別」タスクによる検証．}
    \label{fig:overview}
\end{figure*}

\section{関連研究}
\subsection{画像エンコーダと対照学習による意味理解}
ロボットの環境認識において，視覚情報の解析は最も基本的な課題の一つである．従来はCNNやVision Transformer（ViT）を用いた物体認識が主流であったが，近年では，画像と自然言語を共通の埋め込み空間で結びつける対照学習が大きな進展をもたらしている．

その代表例であるCLIP（Contrastive Language-Image Pre-training)~\cite{CLIP}は，画像とその説明文のペアを用い，対応するペアの類似度を上げ，無関係なペアの類似度を下げるように学習を行う．この手法の利点は，単なるカテゴリ分類に留まらず，画像が持つ「意味」を言語ベースで抽出できる点にある．これにより，ロボット分野においてもCLIPort~\cite{shridhar2021cliport}に代表されるように，未学習の物体や複雑な指示に対しても，言語に基づく外観・幾何学的特徴の意味的理解を介して柔軟に動作を決定するアプローチが可能となった．

しかし，これら視覚ベースの手法が対象とする意味理解は，本質的に物体の外観やジオメトリのドメインに限定されている．
視覚情報からは，容器の内容物の有無などに代表される内部状態や，微細な滑りなどの環境との動的な相互作用といった力学的情報を直接抽出することは困難である．
本研究では言語との意味的統合というアプローチを踏襲しつつも，その対象を「視覚的には識別不可能な力学的ダイナミクスの意味理解」へと拡張する．

\subsection{触覚センシングの展開と力覚情報の優位性}
視覚に依存する手法の限界を補完するため，ロボットの指先に視触覚センサ（GelSight~\cite{GelSight}等）を搭載し，接触面の変形を画像として捉える研究が盛んに行われている．
特に\cite{CLTP, BindingTouch, MMWand}らは，対照学習により物体把持時の触覚画像を対応する視覚画像や言語記述と結びつけることで，材質や質感の予測を可能にした．これらは表面のテクスチャや微細な形状変化の認識には極めて有効であり，近年では大規模言語モデル (LLM) と統合したTactile Language Model~\cite{yu2024octopi}も提案されている．

しかし，視覚をベースとした触覚センシングは
画像処理に起因する計算遅延やサンプリングレートの制約から，初期滑りなどの高周波な物理イベントへの即応性に課題を残している．また，指先の局所的な変形のみを画像という静的なモダリティとして捉えるため，物体を介してアーム全体にかかる環境反力や，物体の慣性・重心といった大域的・ダイナミックな力学特性を計測することができない．

一方，高周波（\qty{100}{\Hz}以上）かつ大域的な計測が可能な6軸力覚センサを用いた深層学習アプローチ（ForceVLA~\cite{ForceVLA}等）も存在する．だが，これらの既存手法はタスク成功率の向上を目的としたEnd-to-Endの行動生成に特化しており，力覚信号そのものが持つ物理的な意味を汎用的な言語空間に射影する試みは見られない．
インピーダンス制御や特定の閾値判定などの力覚制御もまた，タスクや環境に依存したハンドクラフトな設計を必要とするため，汎用的な意味論への接続には至っていない．

本研究の新規性は，これら「視触覚ベースの意味理解」と「力覚ベースの行動生成」の隙間に存在する技術的な隔たりを埋める点にある．
すなわち，6軸力覚センサの持つ高周波・大域的という独自の物理的優位性を最大限に活かしつつ，End-to-Endのタスク特化型学習から脱却し，多様な物理イベントを汎用的な言語空間へとマッピングする「力触覚エンコーダ」を構築する．


\section{提案手法}
\subsection{力触覚と言語のアライメントに基づく問題定式化}
本研究の目的は，6軸の力触覚（Force/Torque, F/T）データの時系列を，対応する自然言語の記述と意味的に整合した特徴量空間へと射影する関数を学習することである．

既存の対照学習アプローチ~\cite{CLIP, CLTP, BindingTouch}は，画像のような高次元かつ静的なデータ同士の整合を前提として設計されている．これに対し本研究が扱う力触覚時系列は，12次元という極めて低次元かつ，時間方向にのみ高周波な動的時系列である．このように「豊かな概念を持つ自然言語と，低次元なセンサ信号という非対称な二つのモダリティをいかにして共通空間で整合させるか」が，本定式化における最大の課題である．

この非対称性を解消するためには，力触覚時系列における時間方向の大域的な動的推移と，衝撃などの局所的で急峻な変動の双方を精緻に捉えて特徴量を抽出することが極めて重要となる．そこで本研究では，この力学的特性を効果的にモデル化すべく，後述するように複数の異なる時系列アーキテクチャを導入し，その表現能力を包括的に探索する．

さらに，すでにセマンティクスが確立されている言語の特徴量空間に対し，力触覚時系列のダイナミクスを歪めることなく組み込むための定式化を行う．具体的には，先行研究~\cite{CLTP, BindingTouch}の知見に倣い，事前学習済みの言語特徴量空間を固定した上で，力触覚エンコーダのみを最適化する．この片方向のアライメントにより，力触覚の時系列変動が，言語空間の持つ抽象的な物理概念へと誘導され，言語の意味空間への統合が可能となる．

時刻 $t$ における力触覚時系列を $\mathbf{X}_t = [\mathbf{x}_t, \cdots, \mathbf{x}_{t+T}]$ とする．ここで，$\mathbf{x}_\tau \in \mathbb{R}^{D_\text{sensor}}$ は時刻 $\tau$ における測定値の集合であり，本検証では左右の6軸力覚センサの測定値なので12次元となる．
我々は，これを言語と整合した表現 $\mathbf{e}^\text{F/T}_t \in \mathbf{D}_{out}$ に写像する関数 $\mathbf{f}^\text{F/T}$ の学習を目指す．
\begin{align*}
&\mathbf{e}^\text{F/T}_t = \mathbf{f}^\text{F/T}(\mathbf{X}_t) \
&\text{s.t.} \quad \mathbf{e}^\text{F/T}_t \sim \mathbf{e}^\text{Lang}_t
\end{align*}
ここで，$\mathbf{e}^\text{Lang}_t \in \mathbf{D}_{out}$ は時刻 $t$ における自然言語記述 $l_t$ の埋め込み表現である．

これを実現するため，対照学習を導入する．
各力触覚時系列 $\mathbf{X}_t$ とそれに対応するテキスト記述 $l_t$ のペアに対し，両方のエンコーダを用いて埋め込みを計算する．
\begin{align}
\mathbf{e}^{\text{F/T}}_t &= \mathbf{f}^{\text{F/T}}(\mathbf{X}_t) \\
\mathbf{e}^{\text{Lang}}_t &= \mathbf{f}^{\text{Lang}}(l_t)
\end{align}
ここで，$\mathbf{f}^{\text{Lang}}$ は事前学習済みのテキストエンコーダであり，学習中はパラメタが固定される．
これらの埋め込みを共通の意味空間で整合させるために，CLIPのトレーニングで一般的に使用されるInfoNCE損失\cite{InfoNCEloss}を採用した．
\begin{align}
    \mathcal{L}_{\text{InfoNCE}} = -\frac{1}{B}\sum_{i=1}^{B} \log \frac{\exp(\text{sim}(\mathbf{e}^{\text{F/T}}_i, \mathbf{e}^{\text{Lang}}_i) / \tau)}{\sum_{j=1}^{N} \exp(\text{sim}(\mathbf{e}^{\text{F/T}}_i, \mathbf{e}^{\text{Lang}}_j) / \tau)}
\end{align}
ここで，$\text{sim}(\cdot, \cdot)$ はコサイン類似度，$B$ はバッチサイズ，$\tau$ は学習可能な温度パラメータである．
この目的関数は，ペアになっている力触覚時系列と言語の埋め込み間の類似度を最大化し，同時に各バッチ内のペアではないサンプル間の類似度を最小化するようにモデルを促す．

このように学習された表現は，潜在的な物理ダイナミクスを効果的に捉えることができ，下流の操作タスクにおいて有用であることが期待される．

\subsection{モデル設計}
力触覚の時空間的な特徴を効果的に抽出するため，本研究では以下の2つの異なるアーキテクチャを検討した．なお，モデル間の公正な比較のため，モデルのパラメータ総数は約1Mに統一している．

\textbf{1. PatchTST (Patch Time Series Transformer)．}
PatchTST\cite{PatchTST}は，時系列処理のための最先端のTransformerベースのモデルである．
ある時間セグメントの力触覚時系列 $\mathbf{X}_t$ は，まず時間軸に沿って長さ $P_t$ の重なりのないパッチに分割される．
各パッチは線形射影によってトークンに変換され，位置埋め込みが加算される．
その後，Transformerブロックが各チャネルを独立に処理し，個々の測定値の時間的なダイナミクスに焦点を当てる．
最後に，バックボーンの出力 $\mathbf{H}_t = [\mathbf{h}_t^1, \cdots, \mathbf{h}_t^{D_\text{sensor}}] \in \mathbb{R}^{D_\text{backbone} \times D_\text{sensor} }$ がMLPによって集約され，センサ間の関係性を捉えることで，ターゲットセグメントの埋め込み表現 $\mathbf{e}^\text{F/T}_t$ を生成する．

\textbf{2. CNN (Convolutional Neural Network)．}
CNNは，時系列データ内の局所的な時間変化（パルス状の衝撃や急激な滑りなど）を特徴量として抽出するために採用した．力触覚時系列の時間軸方向に1次元畳み込み（1D Convolution）を適用し，層を重ねることで低次から高次の物理的特徴を階層的に学習する．畳み込み後の特徴マップはグローバル平均プーリング等によって集約され，最終的な全結合層を経て埋め込み表現 $\mathbf{e}^\text{F/T}_t$ を生成する．PatchTSTが全域的な依存関係を捉えるのに対し，CNNは受容野内でのパターン認識能力に特化した比較対象として位置づける．


\section{データセット構築}
\subsection{データ収集環境とモダリティ}
本研究ではシミュレーション環境において，ロボットアームによるマルチモーダルなデータセットを構築した．シミュレータとしてはGenesis~\cite{Genesis}を用い，ロボットアームはFranka Emika Pandaを採用した．各試行において以下の3種類のデータを同期して収集した．

\begin{itemize}
    \item \textbf{力触覚データ．} 両指のグリッパに仮想的な6軸力覚センサを搭載し，\qty{100}{\Hz} でサンプリングした．解析の際は，連続したシーケンスを 窓長\qty{0.8}{s}（80ステップ），ストライド\qty{0.1}{s}のウィンドウでセグメント化して用いる．
    \item \textbf{画像データ．} アーム手首の一人称視点および正面の三人称視点の2視点からRGB画像をキャプチャした．これらはベースラインとの比較評価および視覚的な状況把握に用いる．
    \item \textbf{言語アノテーション．} シミュレータ内の物理情報に基づき，後述するルールベース手法により自動生成した．
\end{itemize}

\subsection{データ収集タスクの設計}
多様な力学的イベントを網羅するため，ランダムに配置されたタイルへ物体を運ぶ「Pick-and-Place」動作を基本とし，以下の要素を組み合わせてデータを生成した．
アームの基本動作は，把握（grasp），持ち上げ（lift），保持（hold），下降（descend），落下（drop），配置（place），押し付け（press）の7種類で構成される．特に配置動作においては，「タイルの上で離す（落下）」，「接地後即座に離す（配置）」，「一定の負荷をかけてから離す（押し付け）」といったバリエーションを持たせ，接触状態の差異を学習可能にした．

\subsection{ルールベースによる言語アノテーション}
\label{subsec:annotation}
客観的かつ一貫性のある言語記述を生成するため，「Action」「Weight」「Interaction」の3つの物理的側面を抽出し，これらを組み合わせる自動アノテーション手法を採用した．各要素の具体的な判定基準と分類内訳は以下の通りである．

\begin{enumerate}
    \item \textbf{Action．} ロボットの制御ステートおよび突発的な物理イベントに基づき判定する．アームの基本動作である grasp, lift, hold, drop, place, press に，マニピュレーション中に不安定さから意図せず物体を落としたケースである accidental drop を加えた計7クラスに分類される．
    \item \textbf{Weight．} 物体の質量に基づき，\qty{100}{g} を閾値として light または heavy の2クラスに分類する．この閾値は，飲料缶やペットボトル等の家庭内オブジェクトにおいて，中身が空であるか充填されているかを識別するという実用的なユースケースを想定して設定した．
    \item \textbf{Interaction．} 指の座標に対する物体重心の相対速度を算出し，2段階の閾値設定により stable（安定保持），slip slowly（緩やかな滑り），slip quickly（急激な滑り）の3クラスに分類する．滑りの閾値は，物体が把持点から逸脱し始める微細な変動（Stick-Slip現象）を捉えられるよう調整した．
\end{enumerate}

これらを統合し，「Lifting a heavy object, letting it slip quickly.」あるいは不慮の脱落時には「Accidentally dropped a heavy object.」といった，物理現象を詳細に記述する自然言語文（全46種類）を各ウィンドウに対して自動付与した．

\subsection{データセットの統計と品質管理}
学習モデルの偏りを防ぐため，特定の動作が全体の15\%を超えないようサンプリングを調整し，重量についてもlight:heavyが概ね1:1の比率となるようデータバランスを調整した．

最終的なデータセットは，112,695個の学習サンプル（Google Scanned Objects~\cite{GoogleScanObject}を使用）と5,895個の評価サンプル（YCB Objects~\cite{YCB}を使用）で構成されている．未知の物体（YCB）に対する評価を行うことで，モデルの形状汎化性能を検証可能としている．

\section{力触覚エンコーダの評価}
\label{sec:eval}
本章では，対照学習によって得られた力触覚エンコーダが，力触覚時系列を物理状態に対応する言語的意味空間へ適切に写像できているかを評価する．
具体的には，エンコーダから得られる埋め込みと，各物理状態を表すテキスト埋め込みとのコサイン類似度に基づき，物理イベント，物体の重さ，および把持状態の分類性能を定量的に評価する．
さらに，混同行列および力触覚時系列の可視化を用いて，識別が容易な物理状態と困難な物理状態を分析する．
最後に，UMAPによる特徴空間の可視化を通じて，学習された力触覚表現と言語埋め込みの空間的な対応関係を確認する．

\subsection{テキスト埋め込みとの類似度に基づく物理状態分類}
\label{subsec:classification}
学習した力触覚エンコーダから得られる埋め込みと，各ラベルに対応するテキスト埋め込みとのコサイン類似度を算出することで，物理事象の識別性能を評価した．
本評価では，提案する力触覚エンコーダの実装としてCNNおよびPatchTSTを用い，時間的構造を明示的に扱わないベースラインとして，入力時系列 $\mathbf{X}_t$ をFlattenして言語空間へ射影するMLPを比較した．
分類対象は，前節で定義した物理的な3つの側面，すなわち \textbf{action}（7クラス：accidental drop, drop, grasp, hold, lift, place, press），\textbf{weight}（2クラス：light, heavy），\textbf{interaction}（3クラス：slip quickly, slip slowly, stable）である．表~\ref{tab:f1_scores_comparison}に各モデルのMicro F1スコアを示す．

CNNおよびPatchTSTはいずれもMLPを上回る性能を示した．
この結果は，力触覚時系列と言語ラベルの対応付けには，単純な全結合構造よりも，時間的な局所変化や時系列パターンを明示的に扱う構造が有効であることを示している．
また，すべてのモデルにおいて一定以上の分類性能が得られたことから，対照学習により，力触覚時系列を物理状態に対応する言語的意味空間へ射影できていることが確認できる．

CNNはaction，weight，interactionのすべてにおいて最も高いスコアを示し，特にinteraction分類において優位であった．
これは，滑りなどの物理現象が短時間の局所的な力変化として現れやすく，畳み込み構造がそのような局所パターンの抽出に適しているためであると考えられる．
一方，PatchTSTもactionおよびweight分類においてCNNと同程度の性能を示した．
なお，後続のアプリケーション評価で扱う接地判定および重さ判定は，それぞれaction分類とweight分類に対応している．
そのため，本研究では時系列全体の文脈を扱えるPatchTSTを代表モデルとして用いる．

\subsection{混同行列に基づく物理イベントの分析}
\label{subsec:heatmap}
図~\ref{fig:label_heatmap}の混同行列に基づき，物理イベント毎の識別傾向を分析する．
まず，「accidental drop（意図しない脱落）」と「drop（意図的な解放）」が明確に識別されている点は注目に値する．これらは視覚的には類似した「物体が離れる」挙動であるが，指を開く際の反力変化の有無という力学的差異をモデルが正確に捉えていることを示唆している．

一方で，「hold（保持）」と「press（押し付け）」の間には一定の混同が見られた．
これは，いずれも一定の負荷が継続的に加わる静的な力特性を共有しているためであると考えられる．
本研究のモデルは，力の大きさや時間的変化に基づく状態識別には有効である一方で，「どの方向に力が作用しているか」や「力が物体保持に由来するのか，環境への押し付けに由来するのか」といったベクトル的・関係的な意味の抽出には課題が残る．
この点は，物体，グリッパ，環境の関係性をより明示的に扱うモデル設計によって今後改善できる可能性がある．

\subsection{滑り状態の識別と失敗要因の分析}
interaction分類では，急激な滑り（slip quickly）は比較的高い精度で検知できた一方で，緩やかな滑り（slip slowly）の識別は困難であった．
図~\ref{fig:slip_stable_ft}に示すように，滑り発生時にはstick-slip現象に起因する微細なスパイクが力触覚時系列に現れる．
特に急激な滑りでは，このスパイクが明瞭に観測されるため，モデルが滑り状態を識別しやすかったと考えられる．

一方，緩やかな滑りでは，力の変化が小さく，安定把持時のノイズや微小な姿勢変化と区別しにくい．
そのため，現在のサンプリング周波数や特徴抽出の時間分解能では，滑りの開始や進行を十分に捉えきれない場合がある．
この結果は，滑り検知のような高周波かつ微小な物理変化を扱うためには，より高いサンプリング周波数での計測や，高周波成分に特化した特徴抽出手法が有効である可能性を示している．

\begin{table}[t] 
    \centering 
    \caption{各モデルのMicro F1スコア} 
    \label{tab:f1_scores_comparison} 
    \small 
    \begin{tabular}{@{}lcccc@{}}
        \toprule
        Model & action & weight & interaction \\
        \midrule
        PatchTST &  0.7530 & 0.7779 & 0.6959 \\
        CNN & 0.7571 & 0.7849 & 0.7395 \\
        MLP & 0.7374 & 0.7530 & 0.7005 \\
        \bottomrule
    \end{tabular}
\end{table}

\begin{figure}[t]
    \centering
    \includegraphics[width=\linewidth]{fig/label_pred_heatmap.pdf}
    \caption{評価データにおける提案モデル（PatchTSTベース）のAction識別精度．}
    \label{fig:label_heatmap}
\end{figure}

\begin{figure}[t]
    \centering
    \includegraphics[width=\linewidth]{fig/slip_stable_ft.pdf}
    \caption{操作中にグリッパの片指で観測された力触覚時系列データの比較．滑り発生時（左）ではStick-Slip現象に起因する微細なスパイクが確認できるが，安定把持時（右）ではそのような定常的なスパイクは観測刺されていない．}
    \label{fig:slip_stable_ft}
\end{figure}

\subsection{特徴空間の定性的評価}
学習された空間の意味的構造を検証するため，力触覚時系列の埋め込みを7つの基本Actionに従って色分けし，可視化した．さらに，言語との整合性を確認し分析するため，各言語ラベルに対応するテキスト埋め込みベクトルを重ねて描画した．
可視化にはUMAP~\cite{umap}を用いた．まず，評価データセットの力触覚時系列の埋め込みベクトルに対して2次元への変換を計算した．このプロセスにより，ベクトル間の相対的な位置関係を維持しながら次元を削減する．その後，これと同一の変換をテキストベクトルに対しても適用した．

その結果，各ラベルに対応する力覚データが明確なクラスタを形成し，かつ対応する言語ラベルの埋め込みベクトルと近接して配置されていることが確認された．%これにより，対照学習を通じて力学的な時系列データが言語的な「意味」へと適切に射影されていることが示された．
この結果は，対照学習を通じて，力学的な時系列データが「grasp」「hold」「place」などの言語的な物理概念と対応付けられていることを示している．
すなわち，提案手法は力触覚時系列を単なるセンサ特徴としてではなく，言語的に解釈可能な意味空間上の表現として獲得できていると考えられる．

なお，図中の右端において，いずれのテキストラベルとも結びつかない独立したクラスタが確認された．この要因として，本研究の対照学習の設定が挙げられる．本データセットは47種類のテキスト記述に対し数万サンプルの力覚データを対応させているため，バッチサイズ（512）内において同一のテキストラベルを持つ異なる力覚サンプル同士が「ネガティブペア」として扱われ，空間上で遠ざけられた結果，特定の物理的特徴を持つサンプル群が独立したクラスタを形成した可能性がある．これは，ロボットの触覚データのような「ラベル密度の低い」ドメインにおける対照学習特有の現象であり，正解ラベルの重複を考慮した損失関数の導入が今後の解決策として考えられる．

\begin{figure}[t]
    \centering
    \includegraphics[width=\linewidth]{fig/umap_DICOMO.pdf}
    \caption{UMAPを用いた力触覚表現と言語埋め込みの2次元可視化．各ドットは評価データセットの力触覚時系列ベクトル（基本Actionに基づく7色で分類）を示し，×印は対応する自然言語ラベルのテキスト埋め込みベクトルを示す．}
    \label{fig:umap}
\end{figure}

\section{アプリケーション評価}
提案モデルによって学習された力触覚表現が，下流の物体操作タスクに応用可能であることを検証するため，シミュレーション環境下で2種類のアプリケーション評価を行った．
本評価では，接触状態の理解を必要とする「物体の接地を検知して安全に把持を解除するタスク」と，重さの理解を必要とする「外観が類似した物体の中から片付け対象を選別するタスク」を設定した．
いずれのタスクにおいても，タスク固有のファインチューニングや閾値の再設計は行わず，事前学習済みのPatchTSTベースのエンコーダをそのまま用いた．

\subsection{接地状態の理解に基づく安全な物体配置}
ロボットが生活空間で物体を運搬・配置する際には，対象物を目的位置まで移動させるだけでなく，机や棚などの環境と接触した状態を適切に理解し，安全なタイミングで把持を解除する必要がある．
例えば，コップやスマートフォンのような日用品を机上に置く場合，接地前に手を放すと落下を招き，接地後も押し付け続けると物体や環境を損傷する恐れがある．
しかし，視覚情報のみから物体と接触面の微小な距離や接地の瞬間を正確に推定することは難しく，奥行き推定の誤差やオクルージョンの影響を受けやすい．
そこで本タスクでは，力触覚時系列から「把持物体が環境に接地し，荷重が支持面へ移り始めた状態」を意味的に検出し，適切なタイミングでグリッパを開く制御への応用を検証する．

\textbf{実験設定と解放制御．}
対象物体には寸法および密度をランダム化した直方体を用いた．
また，評価時の汎化性能を確認するため，物体の密度および寸法を試行ごとにランダム化した．
タスクは，(1) 接近・把持，(2) 持ち上げ・保持，(3) 下降・接地判定，(4) 解放の順で進行する．
アーム下降時，左右の6軸力覚センサから得られる直近 $0.8\,\text{s}$ 間の時系列データ（12次元）を0.1秒周期で提案エンコーダに入力する．得られた特徴量ベクトルと，テキスト「hold」および「place」の埋め込みベクトルとのコサイン類似度をそれぞれ算出し，「place」との類似度が「hold」を上回った時点で解放動作へ移行する．

\textbf{評価指標とベースライン．}
未知の物体に対する20試行の評価を行い，「接地後解放率（落下や過負荷を起こさず適切に解放できた割合）」および「解放遅延時間（物理的な接地から実際にグリッパが開き始めるまでの時間）」を計測した．なお，物理的な接地から $12\,\text{s}$ が経過しても解放動作が行われない場合は，タイムアウトによる失敗と定義した．

比較対象として，以下の2種類のベースラインを設定した．
\begin{itemize}
    \item[1.] \textbf{$F_z$スパイクベースライン．}力覚のz軸方向（鉛直方向）の変動量（1ステップ前との差分の絶対値）が，あらかじめ設定した閾値を超えた時点を接地と判定し，解放動作に移行する．閾値は，$0.1\,\text{N}$ から $10\,\text{N}$ までの区間を等分した10個の候補に対して，それぞれ6回の予備実験を行い，最も高い成功率を示した値を採用した．%その上で，最終的な20試行の評価を実施した．
    \item[2.] \textbf{視覚ベースライン．}手首の一人称視点および正面の三人称視点の画像を入力とし，事前学習済みVLM（OpenAI gpt-4o）を用いて，物体が床に接地しているかを判定する．シミュレーション環境下では，現実世界と同様に，奥行きや接触面の微細な変化を画像のみから判別することが困難である．そのため，本ベースラインでは，床面へのタイルの設置，輪郭が明瞭な物体モデルの選定，および複数のカメラアングルの調整を行い，最も精度の高かった環境設定を採用して評価した．
\end{itemize}


\textbf{結果と考察．}
提案モデルは 85\% の高い確率で，落下を伴わずに接地後の把持解除に成功した．成功試行における解放遅延時間は平均 $0.63\,\text{s}$，最大 $1.14\,\text{s}$ であった．
一方，$F_z$スパイクベースラインの成功率は50\%に留まった．この手法は遅延時間こそ $0.01\,\text{s}$ と極めて短いものの，物体の質量，姿勢，下降速度に大きく依存するため，タスクに応じた閾値調整が不可欠となる．さらに，滑り（Stick-Slip現象）によるノイズを接地スパイクとして誤検知し，物体が空中にある状態で手放してしまうケースが散見された．また，視覚ベースラインでは，落下を伴わずに物体を解放できた試行は65\%で，残りの試行では接地を検知できず，把持を継続したままタイムアウトした．照明条件による反射やオクルージョン，奥行き情報の不足により，画像のみから接地の瞬間を安定して捉えることは困難であった．

これに対し，提案モデルは直近 $0.8\,\text{s}$ の力触覚時系列から，単一時刻の力の大きさではなく，接近，接触，荷重移動に至る時間的な変化を捉えることができる．そのため，瞬間的なノイズに過度に反応することなく，多様な物体に対してロバストに接地状態を判断し，安全な把持解除を実現できることが示された．

\begin{figure}[t]
    \centering
    \includegraphics[width=\linewidth]{fig/place.pdf}
    % \caption{設置動作時の視覚入力と力触覚時系列．左の画像では物体はタイルに触れていないが，これを視覚的に判断するのは難しい（ただし，シミュレーションの描画限界にも起因する）．一方，物体接地時の力触覚時系列にはその手応えが明確に表れる．}
    \caption{物体設置動作時における視覚入力と力触覚時系列データの対比．左図はアーム下降時（接地はしていない）における一人称視点のRGB画像を示し，右図は接地前後のプロセスにおけるグリッパの力触覚時系列を示す．}
    \label{fig:place}
\end{figure}

\subsection{重さの理解に基づく搬送対象の選別}
提案モデルが物体の重さに関する力触覚的な意味を識別できるかを検証するため，外観が同一の2つの物体から重い方を選択して搬送するタスクを実施した．本タスクは，生活支援における多様な実ユースケースを想定している．例えば，飲食後の片付け時に外観から中身の有無が判別しにくい缶やコップの選別，調味料ボトル等の重量変化に基づく自律的な補充判断，さらにはユーザーへの物品手渡し時に荷重の移行から相手の確実な受領を判定するケースなどが挙げられる．
このような場面では，視覚情報だけでは判断が難しい一方で，把持・持ち上げ時に生じる力触覚応答から，物体の重さや内容物の有無を推定できる可能性がある．

\textbf{実験設定と識別手法．}シミュレーション環境内に，外観と寸法が同一で密度のみが異なる2つの立方体を配置した．一方は $100\,\text{g}$ 未満，もう一方は $100\,\text{g}$ 以上となるように設定した．
ロボットは左右の物体を順に把持・持ち上げ，その間の力触覚時系列データを提案エンコーダへ入力する．得られた埋め込みベクトルと，テキスト「heavy」の埋め込みベクトルとのコサイン類似度を算出し，類似度がより高かった方を重い物体として判定し，目標地点へ搬送する．20試行における選択の正答率を評価指標とした．

\textbf{結果と考察．}提案モデルは，物体を持ち上げる動作（Lift）時の力覚データを用いた場合，85\% の精度で重い物体を正しく選択することに成功した．一方，空中で静止・保持している動作（Hold）時のデータを用いた場合の正答率は70\%であった．これは，第5章の分類タスクにおけるweight識別精度（Lift時: 0.81，Hold時: 0.77）の傾向と一致している．%重量という物理特性は，静的な保持状態よりも，慣性や加速度を含む動的な持ち上げ動作において，より顕著な力学的特徴として現れることをモデルが正しく学習していると言える．

比較対象として，z軸方向の力の絶対値平均を直接比較するベースラインを検証した．その結果，本ベースラインは物体持ち上げ時の時系列に適用した場合95\%，保持時の時系列に適用した場合85\%の正答率を示し，本実験設定では提案手法を上回る結果となった．
本検証においてベースラインが提案手法を上回った要因として，直接相対比較が可能な実験設定が挙げられる．本実験では物体の重心や把持位置，ロボットの動作プロファイルが完全に統一されており，さらにノイズのないシミュレーション環境では，アームの加速度に伴う慣性力が$F_z$に線形に反映されるため，2つの値を単に大小比較すれば極めて高精度に識別可能となるためである．

しかし，実際の生活支援シナリオにおいて，ロボットが「この缶には中身が残っているか？」などを判断する際，比較対象となる空の缶が常に隣に用意されているわけではないため，絶対評価の能力が求められる．ベースライン手法で絶対評価を行う場合、$F_z$に対する閾値設計が不可欠となるが、この値はロボットの姿勢や持ち上げ速度の変化、把持位置や物体の重心などに依存するため、動作条件が変わるたびに閾値の再チューニングが必要となる．

これに対し，提案手法は多様なロボット姿勢および物体における力触覚データをあらかじめ学習しており，動的な時系列データから物体の重量に起因する物理特性を抽出している．そのため，\ref{subsec:classification}節での評価の通り，動作速度や姿勢が変動する環境であっても，個別の閾値調整を必要とせず、単一物体の操作時から重量状態をロバストに判断できる汎用性を有している．

\section{限界と今後の展望}
本研究では，力触覚の時系列情報を言語空間に統合することで，接触状態の意味的理解が可能であることを示したが，実用化に向けてはいくつかの課題も残されている．

\textbf{力覚時系列における方向性・微小変化の表現．}\ref{subsec:heatmap}節で述べた通り，現状のモデルでは「保持（hold）」と「押し付け（press）」のように，力のベクトル方向が異なるものの静的な強度が類似している事象の識別において混同が見られた．今後は，単一のエンコーダによる特徴抽出だけでなく，力の方向性を明示的に考慮したネットワーク構造や，高周波成分を強調する周波数領域での解析を組み合わせることで，より高精度な接触状態の同定を目指す．

\textbf{実機環境への適用．}本研究はシミュレーション環境を用いて構築されたが，現実世界のセンサデータには電気的ノイズや温度ドリフト，そして接触面の摩擦特性の不確実性が含まれる．今後は，シミュレーション上で多様な摩擦係数やセンサノイズを模したドメインランダマイゼーションの導入，あるいは少量の実機データを用いたアダプテーション手法を検討し，実環境におけるロバスト性を検証する．

\textbf{行動生成モデルとの統合．}提案エンコーダが抽出した手応えの意味表現をマルチモーダルなポリシーの入力として活用することで，「滑りそうになったから握力を強める」「重いと感じたからゆっくり運ぶ」といった，物理的な理解に基づいたより高度かつ安全な生活支援タスクの遂行が可能になると考えている．

\section{結論}
本研究では，生活支援ロボットの安全かつ確実な作業遂行を目指し，6軸力覚センサの時系列データを自然言語の埋め込み空間へと射影することで，接触状態を意味的に理解可能なエンコーダモデルを構築した．

評価実験において，提案モデルは未知の物体に対しても物理的なイベントを高い精度で識別できることを示した．特に，視覚情報のみでは判別が困難な滑りの検知，および物体の重量判定において高い性能を発揮した．
また，アプリケーション評価を通じて，提案モデルを活用することで，奥行き情報の欠如を補完する「接地の手応えに合わせた把持の解除」や，外観に依存しない「重量ベースの物体選択」といった高度な判断タスクが，特定の閾値設計なしに実現可能であることを証明した．

本研究で得られた「力学的応答を意味空間で解釈する」というアプローチは，ロボットが物理環境におけるリスクを自律的に判断し，人間と安全に協調するための重要な基盤技術となる．今後は，より複雑かつ動的なマニピュレーションへの拡張を通じて，Physical AIの更なる発展に寄与することを目指す．

\section*{謝辞}
本研究の一部は，JSPS科研費(25H01135)の支援を受けて行われました．ここに謝意を表します．

\bibliography{ref}
\bibliographystyle{unsrt}
\end{document}
